\documentclass[11pt]{article}

\addtolength{\oddsidemargin}{-.5in}%
\addtolength{\evensidemargin}{-1in}%
\addtolength{\textwidth}{1in}%
\addtolength{\textheight}{1.7in}%
\addtolength{\topmargin}{-1in}%

%% Standard packages
\usepackage{amsmath, amssymb}
\usepackage{amsthm}
\usepackage{caption}
\usepackage{mathtools}
\usepackage{bm}
\usepackage{bbm}
\usepackage{graphicx}
\usepackage{booktabs}
\usepackage{hyperref}
\usepackage{cleveref}
\usepackage{algorithm}
\usepackage{algorithmic}
\usepackage{xcolor}
\usepackage{microtype}
\usepackage{natbib}
\providecommand{\theHalgorithm}{\arabic{algorithm}}
\providecommand{\theHfigure}{\arabic{figure}}

\theoremstyle{definition}
\newtheorem{definition}{Definition}[section]
\newtheorem{proposition}[definition]{Proposition}
\newtheorem{remark}[definition]{Remark}

\newcommand{\acks}[1]{%
  \section*{Acknowledgments}%
  #1%
}

%% Notation macros
\newcommand{\R}{\mathbb{R}}
\newcommand{\N}{\mathbb{N}}
\newcommand{\bx}{\mathbf{x}}
\newcommand{\bX}{\mathbf{X}}
\newcommand{\bz}{\mathbf{z}}
\newcommand{\bZ}{\mathbf{Z}}
\newcommand{\bd}{\mathbf{d}}
\newcommand{\bD}{\mathbf{D}}
\newcommand{\bB}{\mathbf{B}}
\newcommand{\bI}{\mathbf{I}}
\newcommand{\bJ}{\mathbf{J}}
\newcommand{\bU}{\mathbf{U}}
\newcommand{\bV}{\mathbf{V}}
\newcommand{\bW}{\mathbf{W}}
\newcommand{\norm}[1]{\left\|#1\right\|}
\newcommand{\abs}[1]{\left|#1\right|}
\newcommand{\innerprod}[2]{\langle #1, #2 \rangle}
\newcommand{\dgeo}{d^{\mathrm{geo}}}
\newcommand{\demb}{d^{\mathrm{emb}}}
\newcommand{\dknn}{d^{k\text{-NN}}}
\newcommand{\Gk}{G_k}
\newcommand{\Stress}{\mathrm{Stress}}
\newcommand{\StressOne}{S_1}
\newcommand{\sgn}{\operatorname{sgn}}
\newcommand{\diag}{\operatorname{diag}}
\newcommand{\placeholderfigurebox}[2][0.9\linewidth]{%
  \fbox{\parbox[c][0.18\textheight][c]{#1}{\centering\footnotesize #2}}%
}
\DeclareMathOperator*{\argmin}{arg\,min}

\begin{document}

\title{Geodesic Multidimensional Scaling:\\
       Self-Consistent Dimensionality Reduction\\ via Graph Geodesic Distances}

\author{Pawel Gajer and Jacques Ravel}
% \author{\name Pawel Gajer \email pgajer@gmail.com \\
% \addr Center for Advanced Microbiome Research and Innovation (CAMRI), Institute for Genome Sciences, Department of Microbiology and Immunology, University of Maryland School of Medicine \\
% Baltimore, Maryland, USA
% \AND
% \name Jacques Ravel \\
% \addr Center for Advanced Microbiome Research and Innovation (CAMRI), Institute for Genome Sciences, Department of Microbiology and Immunology, University of Maryland School of Medicine \\
% Baltimore, Maryland, USA}

\maketitle

%% ============================================================

The central task of \emph{multidimensional scaling} (MDS) is to find a
configuration of $n$ points $\bZ = (\bz_1, \ldots, \bz_n) \in \R^{n \times
  d}$ in a low-dimensional Euclidean space $\R^d$ such that the Euclidean
distances $\norm{\bz_i - \bz_j}$ are as close as possible to a given $n
\times n$ symmetric, non-negative dissimilarity matrix $\Delta = (\delta_{ij})$
\citep{torgerson1952multidimensional, kruskal1964multidimensional,
  borg2005modern, cox2000multidimensional}.  The criterion of closeness is made
precise by a \emph{stress} function, and the embedding is obtained by
minimizing this stress.  MDS and its relatives are among the oldest and most
widely used tools in multivariate statistics, visualization, and machine
learning \citep{shepard1962analysis, sammon1969nonlinear,
  vandermaaten2008tsne, mcinnes2018umap}.

When the input data are given as a matrix $\bX \in \R^{n \times m}$ of $n$
observations in a high-dimensional ambient space $\R^m$, the most common choice
is $\delta_{ij} = \norm{\bx_i - \bx_j}_{\R^m}$, the Euclidean distance in the
original space. This is appropriate when the data are believed to be sampled
uniformly from a convex region of a linear subspace of $\R^m$, so that Euclidean
distances faithfully reflect intrinsic structure. However, when the data has a
complex structure (lies on or near a nonlinear manifold embedded in $\R^m$ or is
even more complex singular space), straight-line Euclidean distances are poor
proxies for the intrinsic distances that govern the geometry of the data: two
points near opposite ends of a curved surface may be close in ambient Euclidean
distance but far apart along the surface.

[A figure showing a 3d embedding of a iknn graph associated with a 16S vaginal
or gut dataset - Illustration of complex data structure showing why Euclidean
distance does not make sense to measure distance between points in $\R^m$ that
are highly structured forming corridors and other complex structures]

This limitation motivated \emph{Isomap} \citep{tenenbaum2000global} that
approximates the intrinsic geodesic distance between points by constructing a
$k$-nearest-neighbor ($k$-NN) graph on the data, computing all-pairs shortest
paths on this graph (weighted by Euclidean edge lengths), and then feeding these
graph-geodesic dissimilarities $\delta_{ij} = \dknn(\bx_i, \bx_j)$ into
classical MDS. For data lying on a smooth, sufficiently sampled Riemannian
manifold, the graph-geodesic distances converge to the true manifold geodesics
as $n \to \infty$ \citep{bernstein2000graph, tenenbaum2000global}, and the
resulting embedding recovers the intrinsic geometry with high fidelity.

Despite this advance, Isomap retains a fundamental asymmetry that has not
previously been identified as a source of systematic distortion.  On the
\emph{input} side, distances are geodesic---shortest paths through the data
graph.  On the \emph{output} side, the stress function measures faithfulness
via \emph{Euclidean} distances between embedded points.  This means that
Isomap optimizes an embedding in which the distance between $\bz_i$ and
$\bz_j$ is the straight-line distance $\norm{\bz_i - \bz_j}$, not the
embedded length of the corresponding input-side geodesic.

The comparison below summarizes the distinction between standard MDS,
Isomap, and the proposed Geodesic MDS in terms of the distance used on
the input and output sides of the embedding problem.

\begin{center}
  \small
  \begin{tabular}{@{}p{0.19\textwidth}p{0.34\textwidth}p{0.34\textwidth}@{}}
    \toprule
    Method & Input dissimilarity $\delta_{ij}$ & Output distance $d_{ij}$ \\
    \midrule
    Standard MDS & Euclidean $\norm{\cdot}_{\R^m}$ & Euclidean $\norm{\cdot}_{\R^d}$ \\
    Isomap & $d_{G_k}$ (graph geodesic) & Euclidean $\norm{\cdot}_{\R^d}$ \\
    Geodesic MDS & $d_{G_k}$ (graph geodesic) & embedded length of a fixed input geodesic \\
    \bottomrule
  \end{tabular}
\end{center}

For data living on a smooth manifold with slowly varying curvature this
asymmetry causes only a small error: nearby embedded points have
approximately equal Euclidean and geodesic distances.  But for data with
genuine graph or tree structure---where the topology enforces that all
paths between certain pairs of points must pass through bottleneck
vertices---the Euclidean output distance is a systematic \emph{underestimate}
of the embedded length of the relevant graph geodesic (by the triangle inequality, with
equality only when points are collinear in the embedding space).

Consider a concrete example: a star graph with center $c$ and three arm
tips $v_1, v_2, v_3$, each arm of unit length.  The graph distance between
$v_1$ and $v_2$ is $\delta(v_1, v_2) = 2$ (through $c$).  Any faithful
embedding must respect the topological constraint that the only path from
$v_1$ to $v_2$ passes through $c$.  Standard MDS and Isomap both try to
place $\bz_{v_1}$ and $\bz_{v_2}$ at Euclidean distance $2$ apart, which
is perfectly achievable in $\R^2$---but doing so forces the embedded arms
to be nearly collinear, misrepresenting the star topology.  The \emph{correct}
criterion is that the \emph{path-length distance along the embedded arms}
from $\bz_{v_1}$ to $\bz_{v_2}$ equals $\norm{\bz_{v_1} - \bz_c} +
\norm{\bz_c - \bz_{v_2}} = 2$, which is compatible with any angular
arrangement of the arms and correctly reflects the star structure.
Figure~\ref{fig:star_bias} is intended to make this contrast visually explicit.

\begin{figure}[t]
  \centering
  \includegraphics[width=0.98\linewidth]{figures/fig_star_bias.pdf}
  \caption{Branching metrics expose the output-side inconsistency.  In a star graph, the graph distance between leaves is realized through the center.  Matching this distance with Euclidean output distances biases the embedding toward collinearity, whereas measuring output fidelity by the embedded lengths of the corresponding input geodesics preserves the correct topological structure.}
  \label{fig:star_bias}
\end{figure}


%% ============================================================
\section{Background}
\label{sec:background}

\subsection{Classical MDS}
\label{sec:classical_mds}

Let $\Delta = (\delta_{ij})$ be an $n \times n$ symmetric dissimilarity
matrix with $\delta_{ii} = 0$.  Classical MDS
\citep{torgerson1952multidimensional, young1938discussion} seeks a
configuration $\bZ \in \R^{n \times d}$ such that
$\norm{\bz_i - \bz_j} \approx \delta_{ij}$.  When $\Delta$ contains
squared Euclidean distances, the solution is obtained by eigendecomposition
of the \emph{double-centered} Gram matrix
\begin{equation}
  \bB = -\tfrac{1}{2} \bJ \bD^{(2)} \bJ,
  \label{eq:gram}
\end{equation}
where $\bD^{(2)}_{ij} = \delta_{ij}^2$ and
$\bJ = \bI_n - n^{-1}\mathbf{1}\mathbf{1}^\top$ is the centering matrix.
The $d$-dimensional embedding is
$\bZ = \bU_d \Lambda_d^{1/2}$, where $\bU_d$ and $\Lambda_d$ contain the
$d$ largest eigenvectors and eigenvalues of $\bB$
\citep{eckart1936approximation, gower1966some}.
This is also known as \emph{principal coordinate analysis} (PCoA)
\citep{gower1966some}.  Classical MDS is a closed-form, non-iterative
method; it yields an exact isometric embedding when $\bB$ is positive
semidefinite (equivalently, when $\Delta$ satisfies Euclidean metric
conditions \citep{schoenberg1935remarks}).

\subsection{Metric MDS and Stress}
\label{sec:metric_mds}

When $\Delta$ is not a Euclidean distance matrix, classical MDS may place
negative eigenvalues in $\bB$, and the embedding is approximate.  Metric
MDS \citep{kruskal1964multidimensional} frames the problem explicitly as
an optimization: find $\bZ \in \R^{n \times d}$ minimizing a \emph{stress}
function that measures discrepancy between embedded distances and input
dissimilarities.

\paragraph{Kruskal's Stress-1.}
The most widely used criterion is Kruskal's \emph{Stress-1}
\citep{kruskal1964multidimensional}:
\begin{equation}
  \StressOne(\bZ) = \sqrt{
    \frac{\displaystyle\sum_{i < j}
      \bigl(d_{ij}(\bZ) - \delta_{ij}\bigr)^2}
    {\displaystyle\sum_{i < j} d_{ij}(\bZ)^2}
  },
  \label{eq:stress1}
\end{equation}
where $d_{ij}(\bZ) = \norm{\bz_i - \bz_j}$ is the Euclidean distance
between embedded points.  The normalization by $\sum_{i<j} d_{ij}^2$
makes $\StressOne$ invariant to uniform scaling of $\bZ$; it takes values
in $[0, 1]$, with 0 indicating a perfect isometric embedding.  An
alternative normalization divides by $\sum_{i<j} \delta_{ij}^2$
\citep{borg2005modern}, and is sometimes called Stress-1r; both are in
common use.

\paragraph{Weighted stress.}
A generalization incorporates pair-specific weights $w_{ij} \geq 0$:
\begin{equation}
  \sigma_w(\bZ)
  = \sum_{i < j} w_{ij}\,\bigl(d_{ij}(\bZ) - \delta_{ij}\bigr)^2.
  \label{eq:weighted_stress}
\end{equation}
Common choices are $w_{ij} = \delta_{ij}^{-2}$ (relative weighting,
emphasizing small distances) and $w_{ij} = 1$ (uniform).  The unweighted
raw stress $\sigma(\bZ) = \sum_{i<j}(d_{ij}(\bZ) - \delta_{ij})^2$
is the numerator of $\StressOne^2$ up to scaling.

\paragraph{Qualitative interpretation.}
Empirical guidelines due to \citet{kruskal1964nonmetric} and
\citet{borg2005modern} suggest that $\StressOne < 0.05$ indicates an
excellent fit, $0.05$--$0.10$ a good fit, $0.10$--$0.20$ a fair fit, and
$> 0.20$ a poor fit, though these thresholds depend on $n$ and the
structure of $\Delta$.

\subsection{SMACOF}
\label{sec:smacof}

Direct minimization of $\sigma$ or $\StressOne$ is a non-convex problem
with many local minima.  The \emph{Scaling by MAjorizing a COmplicated
  Function} (SMACOF) algorithm \citep{deleeuw1977applications,
  deleeuw2009multidimensional} is the dominant approach.  SMACOF exploits
the identity
\begin{equation}
  \sigma(\bZ) = \sum_{i<j}w_{ij}\delta_{ij}^2
  - 2\,\mathrm{tr}(\bZ^\top \bV\,\bZ)
  + \mathrm{tr}(\bZ^\top \bW\,\bZ),
  \label{eq:smacof_decomp}
\end{equation}
where $\bW_{ii} = \sum_{j \neq i} w_{ij}$, $\bW_{ij} = -w_{ij}$ ($i \neq
j$), and $\bV$ is the Guttman transform of $\bZ$:
\begin{equation}
  \bV(\bZ)_{ij} = -w_{ij}\,\frac{\delta_{ij}}{d_{ij}(\bZ)}\quad(i \neq j),
  \qquad
  \bV(\bZ)_{ii} = -\sum_{j \neq i} \bV(\bZ)_{ij}.
  \label{eq:guttman}
\end{equation}
The majorization step replaces the complicated function of $d_{ij}(\bZ)$
by a quadratic upper bound that is tight at the current iterate $\bZ^{(t)}$,
yielding the closed-form Guttman update
\begin{equation}
  \bZ^{(t+1)} = \bW^+ \bV(\bZ^{(t)})\,\bZ^{(t)},
  \label{eq:guttman_update}
\end{equation}
where $\bW^+$ is the Moore-Penrose pseudoinverse of $\bW$.  SMACOF
guarantees monotone decrease of $\sigma$ at each iteration and converges
to a stationary point \citep{deleeuw2009multidimensional}.  Multiple random
restarts are used in practice to mitigate local minima.  Efficient
implementations are available in R \citep{deleeuw2009multidimensional}
and Python \citep{pedregosa2011scikit}.

\subsection{Non-Metric MDS}
\label{sec:nonmetric_mds}

Non-metric MDS \citep{kruskal1964nonmetric, shepard1962analysis} relaxes
the requirement to reproduce input dissimilarity values, requiring only
that the \emph{rank order} of distances be preserved.  The optimization
alternates between isotonic regression (computing optimal
\emph{disparities} $\hat\delta_{ij}$ that are monotone in $\delta_{ij}$)
and gradient steps on the embedding.  Non-metric MDS is appropriate for
ordinal or perceptual data; Geodesic MDS as proposed here is a metric
method, though a non-metric geodesic variant is straightforward to
formulate by applying isotonic regression to graph-geodesic disparities.

\subsection{Isomap}
\label{sec:isomap}

\emph{Isomap} \citep{tenenbaum2000global} addresses the limitation of MDS
on nonlinear manifold-structured data by replacing Euclidean input
dissimilarities with approximations of intrinsic geodesic distances.
Given $\bX \in \R^{n \times m}$, the algorithm proceeds in three steps:

\begin{enumerate}
  \item \textbf{$k$-NN graph construction.}  Build a weighted undirected
        graph $\Gk$ with vertex set $\{1, \ldots, n\}$, connecting $i$ and $j$
        by an edge of weight $\norm{\bx_i - \bx_j}_{\R^m}$ whenever $j$ is
        among the $k$ nearest neighbors of $i$ (or vice versa, in the symmetric
        variant).

  \item \textbf{Geodesic distance computation.}  Compute all-pairs
        shortest-path distances $\dknn(i,j)$ on $\Gk$ using Dijkstra's
        algorithm \citep{dijkstra1959note} or Floyd-Warshall
        \citep{floyd1962algorithm}.  Set $\delta_{ij} := \dknn(i,j)$.

  \item \textbf{Classical MDS.}  Apply classical MDS (Section
        \ref{sec:classical_mds}) to the geodesic dissimilarity matrix $\Delta$
        to obtain the low-dimensional embedding $\bZ \in \R^{n \times d}$.
\end{enumerate}

\paragraph{Convergence theory.}
For data sampled i.i.d.\ from a compact Riemannian manifold $(M, g)$
embedded in $\R^m$, the graph-geodesic distance $\dknn(i,j)$ converges
to the Riemannian geodesic distance $d_M(i,j)$ as $n \to \infty$, provided
$k$ grows appropriately with $n$ \citep{bernstein2000graph,
  tenenbaum2000global, zha2003continuum, donoho2003hessian}.
The resulting classical MDS embedding then recovers the isometric embedding
of $M$ into a Euclidean space (when one exists) \citep{tenenbaum2000global}.

\paragraph{The output-side limitation.}
Despite this theoretical guarantee, Isomap measures the fidelity of the
embedding by Euclidean distances $\norm{\bz_i - \bz_j}$ between embedded
points---not by path-length distances along the embedded manifold.  For a
smooth manifold this is a minor issue: the manifold in the embedding is
nearly flat if $d$ is chosen correctly, so Euclidean and geodesic distances
are approximately equal.  For discrete graph metrics, branching metrics
(trees, stars), or data with topological bottlenecks, the discrepancy is
fundamental, as we elaborate in \cref{sec:geodesic_mds}.

\subsection{The $k$-NN Graph as a Geodesic Structure}
\label{sec:knn_graph}

The $k$-NN graph $\Gk$ constructed in step 1 of Isomap can be understood
as a \emph{metric graph}---a weighted graph whose path metric approximates
some underlying geodesic metric space.  Several observations are relevant
for our construction.

\paragraph{Consistency with the ambient metric.}
The edge weights of $\Gk$ are Euclidean distances in $\R^m$, so the path
metric $d_{\Gk}$ always satisfies $d_{\Gk}(i,j) \geq \norm{\bx_i - \bx_j}$
(the shortest path cannot be shorter than the straight line).  Equality
holds only for pairs connected by a direct edge.

\paragraph{Sparsity and the role of $k$.}
$\Gk$ has at most $kn$ edges (exactly $kn/2$ in expectation for large $n$
with a symmetric construction), so the graph is sparse for small $k$.
As $k$ increases, more edges are added, $d_{\Gk}(i,j)$ decreases toward
$\norm{\bx_i - \bx_j}$, and the graph metric approaches the Euclidean
ambient metric.  At $k = n-1$ (complete graph), $d_{\Gk}(i,j) =
\norm{\bx_i - \bx_j}$ for all pairs, and the construction degenerates to
standard Euclidean MDS.

\paragraph{Spanner interpretation.}
$\Gk$ is a $(1+\varepsilon)$-spanner of the Euclidean metric on $\bX$ for
sufficiently large $k$ depending on the intrinsic dimensionality of the
data \citep{althofer1993sparse, har-peled2006low}.  This means
$d_{\Gk}(i,j) \leq (1+\varepsilon)\norm{\bx_i - \bx_j}$, so the graph
geodesic approximates the Euclidean distance with controlled relative
error.  For data with low intrinsic dimension $d^*$, $k = O(\varepsilon^{-d^*})$
suffices \citep{har-peled2006low}.

\paragraph{Minimum realization.}
An edge $(i,j)$ in $K_n$ (the complete graph on $\bX$) is
\emph{metrically redundant} with respect to $d_{\Gk}$ if there exists a
path through intermediate points with equal total length.  The subgraph
obtained by removing all metrically redundant edges is the
\emph{minimum realization} $G^*$ of the metric $d_{\Gk}$
\citep{dress1984trees, imrich1984optimal}, and encodes the intrinsic
geodesic structure of the data.  For tree metrics, $G^*$ is the unique
spanning tree; for hyperbolic metrics it is sparse; for generic metrics
it approaches the complete graph.
Figure~\ref{fig:k_regimes} is a natural companion here because the role of
$k$ is one of the main geometric ideas of the paper.

\begin{figure}[t]
  \centering
  \includegraphics[width=0.98\linewidth]{figures/fig_k_regimes.pdf}
  \caption{The neighborhood parameter $k$ acts as a geometric scale parameter.  The same horseshoe-shaped point set is shown for three neighborhood sizes, with the same endpoint pair $p,q$ highlighted in each panel.  For small $k$, the graph is sparse and nearly tree-like; for intermediate $k$, shortest paths still follow the manifold and preserve local geodesic structure; and for large $k$, dense shortcut edges appear and the metric approaches the ambient Euclidean regime.}
  \label{fig:k_regimes}
\end{figure}

\subsection{Manifold Learning and Related Methods}
\label{sec:manifold_learning}

The broader context of our work is \emph{manifold learning}: the project
of finding low-dimensional representations of high-dimensional data assumed
to lie near a low-dimensional manifold.  Beyond Isomap, influential methods
include Locally Linear Embedding (LLE) \citep{roweis2000nonlinear},
Laplacian Eigenmaps \citep{belkin2003laplacian}, Hessian Eigenmaps
\citep{donoho2003hessian}, Diffusion Maps \citep{coifman2006diffusion},
$t$-SNE \citep{vandermaaten2008tsne}, and UMAP \citep{mcinnes2018umap}.
With the exception of diffusion maps (which define an intrinsic distance via
Markov chain commute times), all of these methods use Euclidean distances
on the output side, making the inconsistency identified here relevant to the
entire manifold learning literature.

Graph drawing methods---particularly Kamada-Kawai \citep{kamada1989algorithm}
and Fruchterman-Reingold \citep{fruchterman1991graph}---are also related.
Kamada-Kawai minimizes a stress between graph-theoretic shortest-path
distances (in the unweighted graph) and Euclidean distances in the 2D
drawing, which is structurally similar to Isomap applied to an unweighted
graph.  A geodesic variant of Kamada-Kawai, using embedded lengths of
fixed graph geodesics on the output, is a natural companion to Geodesic
MDS but is outside the scope of the present paper.

%% ============================================================
\section{Geodesic MDS}
\label{sec:geodesic_mds}

\subsection{The Self-Consistency Principle}
\label{sec:self_consistency}

We begin by articulating the principle that motivates the proposed method.

\begin{definition}[Self-consistent embedding]
  \label{def:self_consistency}
  Let $(X, \delta)$ be a metric space with a graph structure $G = (X, E,
  w)$ whose path metric $d_G$ satisfies $d_G(i,j) \geq \delta(i,j)$ for
  all $i,j$.  For each pair $(i,j)$, choose a canonical $d_G$-geodesic
  $\gamma_{ij}$ in $G$.  An embedding $f: X \to \R^d$ is
  \emph{self-consistent} with respect to $(G, \Gamma)$, where
  $\Gamma = \{\gamma_{ij}\}_{i<j}$, if the fidelity of the embedding is
  measured by the path length of the image $f(\gamma_{ij})$---i.e., by the
  sum of Euclidean edge lengths along that same input-side geodesic.
\end{definition}

Standard metric MDS and Isomap are not self-consistent in this sense:
they measure fidelity by Euclidean distances $\norm{f(i) - f(j)}$, which
are at most equal to any path-length distance through intermediate vertices.

The self-consistency principle has a natural interpretation:
\emph{the distance used to evaluate the embedding should be the same
  type of distance, along the same family of geodesics, used to construct
  the input dissimilarities.}
If input dissimilarities are graph geodesics, output distances should be
measured along those same graph geodesics after embedding.  This is not
merely a philosophical preference: the
analysis in \cref{sec:systematic_bias} below shows that using Euclidean
output distances introduces a systematic downward bias in the estimated
distances for topologically constrained pairs.

\subsection{Systematic Bias of Euclidean Output Distances}
\label{sec:systematic_bias}

\begin{proposition}[Triangle inequality bias]
  \label{prop:bias}
  Let $\bZ = (\bz_1, \ldots, \bz_n) \in \R^{n \times d}$ be any embedding,
  and let $G = (V, E, w)$ be a weighted graph with vertex set $V = \{1,
  \ldots, n\}$.  For any pair $(i,j)$ with chosen shortest path
  $i = v_0, v_1, \ldots, v_\ell = j$ in $G$, define the embedded length of
  that geodesic by
  \begin{equation}
    \demb_G(i,j;\bZ)
    = \sum_{s=0}^{\ell-1} \norm{\bz_{v_s} - \bz_{v_{s+1}}}.
    \label{eq:emb_geo}
  \end{equation}
  Then
  \begin{equation}
    \norm{\bz_i - \bz_j} \;\leq\; \demb_G(i,j;\bZ),
    \label{eq:triangle_bias}
  \end{equation}
  with equality if and only if all embedded path vertices
  $\bz_{v_0}, \ldots, \bz_{v_\ell}$ are collinear in $\R^d$ in the order
  of the path.
\end{proposition}

\begin{proof}
  By the triangle inequality applied $\ell$ times:
  $\norm{\bz_i - \bz_j} = \norm{\sum_{s} (\bz_{v_{s+1}} - \bz_{v_s})}
  \leq \sum_s \norm{\bz_{v_{s+1}} - \bz_{v_s}} = \demb_G(i,j;\bZ)$.
  Equality holds iff all difference vectors $\bz_{v_{s+1}} - \bz_{v_s}$
  are non-negative scalar multiples of a single unit vector, i.e., the
  points are collinear in the stated order.
\end{proof}


\Cref{prop:bias} shows that standard MDS and Isomap systematically
underestimate embedded geodesic lengths for all topologically
non-degenerate pairs (those whose shortest path has $\ell \geq 2$ and
does not embed as a straight line).  When the input dissimilarity
$\delta_{ij}$ is a graph geodesic, minimizing Euclidean-output stress
creates pressure to satisfy $\norm{\bz_i - \bz_j} \approx \delta_{ij}$,
which in turn forces embedded paths to be approximately collinear.
For tree-structured data this collapses the tree into a lower-dimensional
structure than warranted; for star graphs it forces nearly collinear arm
arrangement regardless of the dimension $d$.

\subsection{The Geodesic Stress Function}
\label{sec:geodesic_stress}

We now define the proposed method.

\paragraph{Setup.}
Let $\bX \in \R^{n \times m}$ be the input data matrix.  Fix a neighborhood
parameter $k \in \{1, \ldots, n-1\}$ and construct the symmetric $k$-NN
graph $\Gk = (V, E_k)$ with $V = \{1, \ldots, n\}$ and
\begin{equation}
  E_k = \{(i,j) : j \in \mathcal{N}_k(i) \text{ or } i \in \mathcal{N}_k(j)\},
  \label{eq:knn_edges}
\end{equation}
where $\mathcal{N}_k(i)$ denotes the $k$ nearest neighbors of $i$ in $\R^m$.
Assign edge weights $w_{ij} = \norm{\bx_i - \bx_j}_{\R^m}$ for $(i,j)
\in E_k$.  Compute the all-pairs shortest-path distance matrix
$\Delta^{(k)} = (\delta^{(k)}_{ij})$ on $\Gk$ using Dijkstra's algorithm.
For each pair $i < j$, choose a canonical shortest path
\begin{equation}
  \gamma_{ij}^{(k)} = (v_0 = i, v_1, \ldots, v_L = j)
\end{equation}
in $\Gk$ with respect to the input edge weights; if several shortest
paths exist, break ties deterministically.  Denote the resulting family
of input geodesics by $\Gamma_k = \{\gamma_{ij}^{(k)} : i < j\}$.

An embedding is a configuration matrix $\bZ = (\bz_1, \ldots, \bz_n)^\top
\in \R^{n \times d}$.  For a fixed graph topology $\Gk$ and fixed
geodesic family $\Gamma_k$, the
\emph{embedded graph} $\Gk(\bZ)$ is the same graph $\Gk$ with edge lengths
replaced by embedded Euclidean distances: edge $(i,j)$ has length
$\norm{\bz_i - \bz_j}_{\R^d}$.  The \emph{embedded length of the chosen
input geodesic} is
\begin{equation}
  \demb_{k}(i,j;\bZ)
  = \sum_{s=0}^{L-1} \norm{\bz_{v_s} - \bz_{v_{s+1}}},
  \label{eq:emb_geodesic}
\end{equation}
where $\gamma_{ij}^{(k)} = (v_0, \ldots, v_L)$.

\begin{definition}[Geodesic stress]
  \label{def:geodesic_stress}
  The \emph{geodesic Stress-1} of an embedding $\bZ \in \R^{n \times d}$
  with respect to the $k$-NN graph $\Gk$ is
  \begin{equation}
    \StressOne^{\mathrm{geo}}(\bZ; k)
    = \sqrt{
        \frac{\displaystyle\sum_{i < j}
              \bigl(\demb_{k}(i,j;\bZ) - \delta^{(k)}_{ij}\bigr)^2}
             {\displaystyle\sum_{i < j}
              \bigl(\delta^{(k)}_{ij}\bigr)^2}
      }.
    \label{eq:geo_stress}
  \end{equation}
  % The \emph{Geodesic MDS} embedding is
  % \begin{equation}
  %   \hat{\bZ}(k) = \argmin_{\bZ \in \R^{n \times d}}
  %   \StressOne^{\mathrm{geo}}(\bZ; k).
  %   \label{eq:geo_mds}
  % \end{equation}
\end{definition}

\begin{remark}[Normalization]
  We normalize by $\sum_{i<j} (\delta^{(k)}_{ij})^2$ rather than
  $\sum_{i<j} (\demb_k)^2$.  This makes $\StressOne^{\mathrm{geo}}$
  independent of the scale of the input dissimilarities and simplifies
  the gradient derivation.  The alternative normalization by
  $\sum_{i<j} (\demb_k)^2$ would introduce a dependency of the
  denominator on $\bZ$, analogous to Kruskal's original Stress-1
  formulation; both choices are valid.
\end{remark}

\begin{remark}[Fixed topology and fixed geodesics]
  The graph topology $E_k$ and the chosen geodesic family $\Gamma_k$ are
  fixed by the input data and are not part of the optimization.  Only the
  vertex positions $\bz_1, \ldots, \bz_n$ are free.  This is a key design
  choice: it ensures that the relevant routes are determined by the
  internal geometry of the input metric space and that the optimization
  preserves that geodesic structure after embedding.
\end{remark}
Figure~\ref{fig:geodesic_stress_definition} should make the fixed-geodesic,
two-metric comparison in \cref{eq:emb_geodesic,eq:geo_stress} easy to see.

\begin{center}
  \includegraphics[width=0.98\linewidth]{figures/fig_geodesic_stress_definition.pdf}
  \captionof{figure}{Definition of geodesic stress.  Left: the input dissimilarity $\delta^{(k)}_{ij}$ is the shortest-path distance between $i$ and $j$ on the fixed $k$-NN graph in the input space.  Right: the output quantity $d_k^{\mathrm{emb}}(i,j;\bZ)$ is the length of the canonical shortest path chosen in the input graph, measured after the vertex positions move in $\R^d$.  The route is fixed by the input geometry; only the embedded edge lengths change during optimization.}
  \label{fig:geodesic_stress_definition}
\end{center}


Note that for a given graph the geodesic Stress-1 function may have infinite
number of global minima. For example, a path graph $v_{0} - v_{1} - v_{2}$ with
any edge lengths, any embedding of that graph into $\mathbb{R}^{2}$ that
preserves edge lengths is a minimum of the geodesic Stress-1 function. An
infinite family of such embeddings can be formed by fixing the edge
$[v_{0},v_{1}]$ embedding and rotating $[v_{1},v_{2}]$ around $v_{1}$.

Let $G$ be a graph with four vertices such that $v_{0}, v_{1}, v_{2}$ form a
triangle and $v_{1}, v_{2}, v_{3}$ form a triangle and the edge $[v_{1},v_{2}]$
is the common edge of these two triangles. Then there are two global minima of
the geodesic Stress-1 function of that graph: one corresponding to vertices
$v_{0}$ and $v_{3}$ situated on the opposite sides of $[v_{1},v_{2}]$ and
another one with $v_{0}$ and $v_{3}$ being on the same side of $[v_{1},v_{2}]$
with the two triangles overlaping each other.

Note that any embedding of the graph that preserves the edge lenghts is a global
minimum of the geodesic Stress-1 function. In general, it may be impossible to
create an embedding where all edge lengths are preserved, so we can aim to
minimize the edge length functional
\[
  E(\bZ) = \sum_{(i,j) \in E} k_{ij} \bigl(\demb(i,j;\bZ) - \delta_{ij}\bigr)^2
\]
where $\delta_{ij}$ is the geodesic distance between $i$ and $j$,
$\demb(i,j;\bZ)$ is defined as above and $k_{ij}$ is a tuning factor. If we want
most of the edges to have their edge lengths preserved we can set $k_{ij}$ to be
the value $dl_{ij}$ of the edge length density $dl$ (possibly after some
normalization of that density) at $\delta_{ij}$ giving us
\[
  E(\bZ) = \sum_{(i,j) \in E} dl_{ij} \bigl(\demb(i,j;\bZ) - \delta_{ij}\bigr)^2
\]

Since the $E$ functional may have many local and global minima, we need another
condition. It is natural to impose a condition that would make the embedding to
be as dispersed as possible in the ambient space in some sense maximizing the
entropy of the embedding. One possible way to specify the maximization of the
spread of the graph in the ambient space is to require that the total repulsion
'force' between all vertices of the graph is minimized by the embedding. Let's
suppose $f_{r}(d)$ is the magnitude of a repulsive force between two vertices
with whose Euclidean distance in the ambient space $l_{ij}$ is $d$. We can
assume that $f_{r}(d)$ is a decreasing function of $d$ and we can define a total
repulsive force
\[
  R(\bZ) = \sum_{i < j} f_{r}(l_{ij})
\]
An embedding $\bZ$ is called a \emph{geodesic MDS embedding} if $\bZ$ is a local
minimum of both $E$ and $R$.

%% ============================================================
%% BIBLIOGRAPHY
%% ============================================================

\acks{Research reported in this publication was supported by grant number INV-048956 from the Gates Foundation.}

\bibliographystyle{plainnat}
\bibliography{geodesic_mds}

\end{document}
